% GENERATED FILE. Edit canonical lessons, then run npm run generate:books.
% canonical-source-hash: fnv1a64:9150f87b88c75d3c
% canonical-lessons: MW-C238-bovno, MW-C238-khodno, MW-C238-ropno, MW-C238-hali, MW-C238-bij

\chapter{Sowing}
\label{ch:mw-sowing}
\hlchaptermodality{238}

\begin{chapteropening}
\textbf{By the end of this chapter:} \emph{I can say I will dig the field, sow millet and plant a tree.}

Complete the last lesson of chapter 238: I can say I will dig the field, sow millet and plant a tree.
\end{chapteropening}

\section[\texorpdfstring{\mw{बोवणो} (bovṇo)}{बोवणो (bovṇo)}]{\mw{बोवणो} (bovṇo) — to sow}

\label{lesson:MW-C238-bovno}



\begin{quote}
Before the new one: say the Marwadi for an egg, then the Marwadi for a frog.
\end{quote}

\subsection*{You'll want to know: \mw{बोवणो}}
\textbf{\mw{बोवणो}} — \emph{bovṇo} — \textquotedblleft{}to sow\textquotedblright{}.

\begin{quote}
\mw{हूं} \mw{काल} \mw{खेत} \mw{मांय} \mw{बाजरो} \mw{बोसूं।}
\end{quote}

\begin{quote}
\emph{hū̃ kāl khet mā̃y bājro bosū̃.}
\end{quote}

\begin{quote}
Tomorrow I'll sow millet in the field.
\end{quote}

Past (I …): \mw{म्हैं} \mw{बाजरो} \mw{बोयो}, \emph{mhaĩ bājro boyo}.

\begin{cousinweb}
\textbf{\mw{बोवणो}} (\emph{bovṇo}) is inherited from Sanskrit \emph{vapati}, \textquotedblleft{}scatters, sows\textquotedblright{}; CDIAL lists Hindi \emph{boyā}, \textquotedblleft{}sows\textquotedblright{}, and names a Gujarati verb beside it. CDIAL gives no Marwari form in this entry (CDIAL 11282).
\end{cousinweb}

\begin{grammarlens}[title={What you've built}]
One more word for this chapter.
\end{grammarlens}

\subsection*{Your turn}
\begin{itemize}
\raggedright
  \item \emph{Say it:} \emph{bovṇo}
  \item \emph{Say it:} \emph{bovṇo}, once more
  \item \emph{Say it:} \emph{hū̃ kāl khet mā̃y bājro bosū̃.} — Tomorrow I'll sow millet in the field.
  \item \emph{Recall:} say the Marwadi for an egg, then the Marwadi for a frog, then say \emph{bovṇo} again
\end{itemize}

\begin{tcolorbox}[breakable,colback=teal!4,colframe=teal!35!black,title={Before you move on}]
What does \mw{बोवणो} mean? (\textquotedblleft{}To sow\textquotedblright{}.) Say it once more.
\end{tcolorbox}
\section[\texorpdfstring{\mw{खोदणो} (khodṇo)}{खोदणो (khodṇo)}]{\mw{खोदणो} (khodṇo) — to dig}

\label{lesson:MW-C238-khodno}



\begin{quote}
Before the new one: say the Marwadi for a frog, then the Marwadi for to sow.
\end{quote}

\subsection*{You'll want to know: \mw{खोदणो}}
\textbf{\mw{खोदणो}} — \emph{khodṇo} — \textquotedblleft{}to dig\textquotedblright{}.

\begin{quote}
\mw{हूं} \mw{खेत} \mw{खोदूं} \mw{हूं।}
\end{quote}

\begin{quote}
\emph{hū̃ khet khodū̃ hū̃.}
\end{quote}

\begin{quote}
I dig the field.
\end{quote}

Past (I …): \mw{म्हैं} \mw{खेत} \mw{खोद्यो}, \emph{mhaĩ khet khodyo}.

\begin{cousinweb}
\textbf{\mw{खोदणो}} (\emph{khodṇo}) is inherited from a reconstructed \emph{*khōdd}, \textquotedblleft{}dig\textquotedblright{}; CDIAL lists Marwari \emph{khodṇo}, \textquotedblleft{}to dig\textquotedblright{} (CDIAL 3934).
\end{cousinweb}

\begin{grammarlens}[title={What you've built}]
One more word for this chapter.
\end{grammarlens}

\subsection*{Your turn}
\begin{itemize}
\raggedright
  \item \emph{Say it:} \emph{khodṇo}
  \item \emph{Say it:} \emph{khodṇo}, once more
  \item \emph{Say it:} \emph{hū̃ khet khodū̃ hū̃.} — I dig the field.
  \item \emph{Recall:} say the Marwadi for a frog, then the Marwadi for to sow, then say \emph{khodṇo} again
\end{itemize}

\begin{tcolorbox}[breakable,colback=teal!4,colframe=teal!35!black,title={Before you move on}]
What does \mw{खोदणो} mean? (\textquotedblleft{}To dig\textquotedblright{}.) Say it once more.
\end{tcolorbox}
\section[\texorpdfstring{\mw{रोपणो} (ropṇo)}{रोपणो (ropṇo)}]{\mw{रोपणो} (ropṇo) — to plant (a seedling or a tree)}

\label{lesson:MW-C238-ropno}



\begin{quote}
Before the new one: say the Marwadi for to sow, then the Marwadi for to dig.
\end{quote}

\subsection*{You'll want to know: \mw{रोपणो}}
\textbf{\mw{रोपणो}} — \emph{ropṇo} — \textquotedblleft{}to plant (a seedling or a tree)\textquotedblright{}.

\begin{quote}
\mw{हूं} \mw{काल} \mw{खेत} \mw{कनै} \mw{रूंख} \mw{रोपसूं।}
\end{quote}

\begin{quote}
\emph{hū̃ kāl khet kanai rū̃kh ropsū̃.}
\end{quote}

\begin{quote}
Tomorrow I'll plant a tree by the field.
\end{quote}

Past (I …): \mw{म्हैं} \mw{रूंख} \mw{रोप्यो}, \emph{mhaĩ rū̃kh ropyo}.

\begin{cousinweb}
\textbf{\mw{रोपणो}} (\emph{ropṇo}) is inherited from a reconstructed \emph{*rupati}, \textquotedblleft{}makes a hole, plants\textquotedblright{}, beside Sanskrit \emph{ropayati}, \textquotedblleft{}plants\textquotedblright{}; CDIAL lists the Old Marwari verb \emph{ropaï}, \textquotedblleft{}plants\textquotedblright{} (CDIAL 10783).
\end{cousinweb}

\begin{grammarlens}[title={What you've built}]
One more word for this chapter.
\end{grammarlens}

\subsection*{Your turn}
\begin{itemize}
\raggedright
  \item \emph{Say it:} \emph{ropṇo}
  \item \emph{Say it:} \emph{ropṇo}, once more
  \item \emph{Say it:} \emph{hū̃ kāl khet kanai rū̃kh ropsū̃.} — Tomorrow I'll plant a tree by the field.
  \item \emph{Recall:} say the Marwadi for to sow, then the Marwadi for to dig, then say \emph{ropṇo} again
\end{itemize}

\begin{tcolorbox}[breakable,colback=teal!4,colframe=teal!35!black,title={Before you move on}]
What does \mw{रोपणो} mean? (\textquotedblleft{}To plant (a seedling or a tree)\textquotedblright{}.) Say it once more.
\end{tcolorbox}
\section[\texorpdfstring{\mw{हाळी} (hāḷī)}{हाळी (hāḷī)}]{\mw{हाळी} (hāḷī) — a ploughman}

\label{lesson:MW-C238-hali}



\begin{quote}
Before the new one: say the Marwadi for to dig, then the Marwadi for to plant (a seedling or a tree).
\end{quote}

\subsection*{You'll want to know: \mw{हाळी}}
\textbf{\mw{हाळी}} — \emph{hāḷī} — \textquotedblleft{}a ploughman\textquotedblright{}.

\begin{quote}
\mw{हाळी} \mw{तड़कै} \mw{खेत} \mw{जावै} \mw{है।}
\end{quote}

\begin{quote}
\emph{hāḷī taṛkai khet jāvai hai.}
\end{quote}

\begin{quote}
The ploughman goes to the field at daybreak.
\end{quote}

A masculine noun.

\begin{cousinweb}
\textbf{\mw{हाळी}} (\emph{hāḷī}) is inherited from Sanskrit \emph{hālika}, \textquotedblleft{}a ploughman\textquotedblright{}, from \emph{hala}, \textquotedblleft{}a plough\textquotedblright{}; CDIAL lists Hindi \emph{hālī} and Gujarati \emph{hāḷi}, \textquotedblleft{}a ploughman\textquotedblright{}. CDIAL gives no Marwari form in this entry (CDIAL 14067).
\end{cousinweb}

\begin{grammarlens}[title={What you've built}]
One more word for this chapter.
\end{grammarlens}

\subsection*{Your turn}
\begin{itemize}
\raggedright
  \item \emph{Say it:} \emph{hāḷī}
  \item \emph{Say it:} \emph{hāḷī}, once more
  \item \emph{Say it:} \emph{hāḷī taṛkai khet jāvai hai.} — The ploughman goes to the field at daybreak.
  \item \emph{Recall:} say the Marwadi for to dig, then the Marwadi for to plant (a seedling or a tree), then say \emph{hāḷī} again
\end{itemize}

\begin{tcolorbox}[breakable,colback=teal!4,colframe=teal!35!black,title={Before you move on}]
What does \mw{हाळी} mean? (\textquotedblleft{}A ploughman\textquotedblright{}.) Say it once more.
\end{tcolorbox}
\section[\texorpdfstring{\mw{बीज} (bīj)}{बीज (bīj)}]{\mw{बीज} (bīj) — a seed}

\label{lesson:MW-C238-bij}



\begin{quote}
Before the new one: say the Marwadi for to plant (a seedling or a tree), then the Marwadi for a ploughman.
\end{quote}

\subsection*{You'll want to know: \mw{बीज}}
\textbf{\mw{बीज}} — \emph{bīj} — \textquotedblleft{}a seed\textquotedblright{}.

\begin{quote}
\mw{म्हारो} \mw{बीज} \mw{आछो} \mw{है।}
\end{quote}

\begin{quote}
\emph{mhāro bīj āchho hai.}
\end{quote}

\begin{quote}
My seed is good.
\end{quote}

A masculine noun.

\begin{cousinweb}
\textbf{\mw{बीज}} (\emph{bīj}) is a tatsama, taken over from Sanskrit \emph{bīja}, \textquotedblleft{}a seed\textquotedblright{} (Monier-Williams); the inherited forms CDIAL lists, such as Hindi \emph{biyā}, have lost the \emph{j} (CDIAL 9250).
\end{cousinweb}

\begin{grammarlens}[title={What you've built}]
One more word for this chapter.
\end{grammarlens}

\subsection*{Your turn}
\begin{itemize}
\raggedright
  \item \emph{Say it:} \emph{bīj}
  \item \emph{Say it:} \emph{bīj}, once more
  \item \emph{Say it:} \emph{mhāro bīj āchho hai.} — My seed is good.
  \item \emph{Recall:} say the Marwadi for to plant (a seedling or a tree), then the Marwadi for a ploughman, then say \emph{bīj} again
\end{itemize}

\begin{tcolorbox}[breakable,colback=teal!4,colframe=teal!35!black,title={Before you move on}]
What does \mw{बीज} mean? (\textquotedblleft{}A seed\textquotedblright{}.) Say it once more.
\end{tcolorbox}
